\documentclass[12pt]{article}
\usepackage[utf8]{inputenc}
\usepackage{amsmath,amssymb,amsthm}
\usepackage{graphicx}
\usepackage{listings}
\usepackage{xcolor}
\usepackage{hyperref}
\usepackage{geometry}
\geometry{margin=1in}
\usepackage{float}

\title{Domácí úkol 10}
\author{Jáchym Kouba}
\date{\today}

% Remove section numbering
\setcounter{secnumdepth}{0}

% MATLAB listings setup
\lstdefinestyle{mystyle}{
  language=Matlab,
  basicstyle=\ttfamily\small,
  keywordstyle=\color{blue},
  commentstyle=\color{gray},
  stringstyle=\color{teal},
  numbers=left,
  numberstyle=\tiny\color{gray},
  stepnumber=1,
  numbersep=8pt,
  tabsize=2,
  breaklines=true,
  showstringspaces=false
}
\lstset{style=mystyle}

\begin{document}
\maketitle

\newpage
\section{Soustava}
Mé datum narození 1.9. 2004, z čehož dostávám $a=1$ a $b=9$. Poslední konstantu
vypočítám jako 

\begin{equation}
  c = 4* \left( \frac{a+b-2}{16}+1 \right) = 6.
\end{equation}

Pomocí těchto konstant dostávám přenos 

\begin{equation}
  \frac{s+9}{(s+6)(s-1)}.
\end{equation}

\section{Příklad}
Mám navrhnout diskrétní regulátor se vzorkovací periodou $T_s = 0.05
\mathrm{s}$, který bude mít následující saturaci pro vstupní signál: $-1.5 \leq
u \leq 1.5$. Musím dále zajistit nulovou ustálenou odchylku při odezvě na vstup,
maximální překmit $30 \%$ a dobu ustálení odezvy na skok nejvýše $3 \mathrm{s}$.
Budu navrhovat PI regulátor, takže implementuji i anti-windup.

Implementace regulátoru bude následující:

\begin{verbatim}
e(k) = r(k) - y(k); % PI controller input
ssum(k + 1) = ssum(k) + e(k); % sum (discrete "integrator")
  
u(k) = kp*e(k) + ki*ssum(k+1);  
\end{verbatim}

A provedení anti-windupu takovéto:

\begin{verbatim}
if abs(u(k)) >= u_sat % anti-windup implementation
    u(k) = kp*e(k);
    ssum(k+1) = 0;
end

% Saturation of the input
u(k) = max(min(u(k), u_sat), -u_sat);
\end{verbatim}

Odladím konstanty $K_p$ a $K_i$ na hodnoty $K_p = 15$ a $K_i = 3$, čímž získám překmit $8.9 \%$ a 
dobu ustálení odezvy na skok $0.7 \mathrm{s}$, což jsou výsledky výrazně lepší, než požaduje zadání.

\begin{figure}[H]
  \includegraphics[width=1\textwidth]{graphs.png}
  \caption{Grafy odezvy systému na skok a výstupu PI regulátoru}
\end{figure}

\end{document}